\documentclass[border=20pt]{standalone}
\usepackage{chemfig}
\begin{document}
% NaH2PO4 : colonne 1 Na, H, H ; colonne 2 trois O ; P central ; colonne 4 un O (coordination)
\chemfig{
    \charge{90:3pt[anchor=south]=$5\delta+$}{P}
    (-[3,2]\charge{90=\|,180=\|,270=\|,135:1pt[anchor=south east]=$\ominus$,45:1pt[anchor=south west]=$\delta-$}{O}-[4,1.5,,,draw=none]\charge{0:1pt[anchor=west]=$\oplus$}{Na})
    (-[4,1.41]\charge{90=\|,270=\|,135:1pt[anchor=south east]=$\delta-$,45:1pt[anchor=south west]=$\delta-$}{O}-[4,1.5]\charge{45:1pt[anchor=south west]=$\delta+$}{H})
    (-[5,2]\charge{90=\|,270=\|,135:1pt[anchor=south east]=$\delta-$,315:1pt[anchor=north west]=$\delta-$}{O}-[4,1.5]\charge{45:1pt[anchor=south west]=$\delta+$}{H})
    (-[0,1.5,,,->]\charge{45=\|,315=\|,0=\|,135:1pt[anchor=south east]=$2\delta-$}{O})
}
\end{document}
